% Deliverable for: 2026-09-04_205253_derive-fase1-capa-delgada
\documentclass[11pt,a4paper]{article}
\usepackage[utf8]{inputenc}
\usepackage[T1]{fontenc}
\usepackage[spanish,es-nodecimaldot]{babel}
\usepackage{amsmath,amssymb,booktabs,array,geometry,microtype}
\usepackage[hidelinks]{hyperref}
\geometry{margin=2.5cm}
\newcommand{\src}[1]{\textsuperscript{\texttt{(fuente: \detokenize{#1})}}}
\newcommand{\dd}{\mathrm{d}}
\newcommand{\ii}{\mathrm{i}}
\newcommand{\avg}[1]{\overline{#1}}
\newcommand{\Rey}{\mathrm{Re}}

\title{Capa viscosa delgada con fondo no deslizante y superficie libre:\\
estructura vertical, fricción lineal y extensión Laplace DN}
\author{Derivación y comprobación reproducible de la Fase 1}
\date{4 de septiembre de 2026}

\begin{document}
\maketitle

\begin{abstract}
Se resuelve el problema vertical de Stokes para dos paredes no deslizantes (DD) y
para fondo no deslizante con tope libre de tensiones (DN). El espectro DN es
\(\lambda_m=\nu[(m+1/2)\pi/h]^2\), por lo que el modo lento produce
\(\alpha=\pi^2\nu/(4h^2)\), cuatro veces menor que en DD. Para ese modo, una
medición PIV en la superficie da \(u(h)/\avg u=\pi/2\). También se deriva la
rama Dirichlet--Neumann que falta en \texttt{laplace\_z} de SPECTER usando sólo
exponenciales acotadas. Estos resultados son exactos para el problema lineal;
su uso como clausura bidimensional requiere las desigualdades de escala y las
salvedades de las secciones~\ref{sec:validez} y~\ref{sec:salvedades}.
\src{resultado_central}
\end{abstract}

\section{Convenciones y alcance}

El fluido es newtoniano e incompresible, tiene viscosidad cinemática constante
\(\nu>0\), ocupa \(0<z<h\), y \(z=0\) es siempre el fondo. ``Libre'' quiere decir
superficie plana e impermeable, \(w(h)=0\), sin tensión tangencial:
\(\partial_z u(h)=\partial_z v(h)=0\). En el fondo,
\(u(0)=v(0)=w(0)=0\). No se modela la deformación de la superficie.
\src{convenciones_fisicas}

Una barra denota promedio vertical,
\begin{equation}
  \avg q(x,y,t)=\frac1h\int_0^h q(x,y,z,t)\,\dd z .
  \label{eq:promedio}
\end{equation}
Las derivadas Neumann de la sección~\ref{sec:laplace} están orientadas en la
dirección \(+z\), que coincide con la variable derivada que construye SPECTER.
\src{promedio_convencion}

\section{Problema vertical de Stokes}
\label{sec:stokes}

Cada componente horizontal satisface el mismo problema escalar
\begin{equation}
  \partial_t u=\nu\,\partial_{zz}u,\qquad 0<z<h.
  \label{eq:calor}
\end{equation}
Buscando \(u(z,t)=Z(z)T(t)\) se obtiene
\begin{equation}
  T(t)=T(0)e^{-\lambda t},\qquad Z''+q^2Z=0,\qquad
  \lambda=\nu q^2>0.
  \label{eq:separacion}
\end{equation}
No hay modo de tasa cero porque el valor de \(u\) está fijado a cero en el
fondo. La misma derivación vale componente a componente para \(v\).
\src{derivacion_stokes}

\subsection{No deslizante--no deslizante (DD)}

Con \(Z(0)=Z(h)=0\), la condición secular es \(\sin(qh)=0\). El espectro
\emph{completo}, los modos y la expansión de un dato inicial \(u_0\in L^2(0,h)\)
son
\begin{align}
 q_n&=\frac{n\pi}{h},&
 Z_n(z)&=\sin\!\left(\frac{n\pi z}{h}\right),\label{eq:ddmodes}\\
 \lambda_n^{DD}&=\nu\left(\frac{n\pi}{h}\right)^2,
 & n&=1,2,\ldots,\label{eq:ddspec}\\
 u(z,t)&=\sum_{n=1}^{\infty}a_n Z_n(z)e^{-\lambda_n^{DD}t},&
 a_n&=\frac2h\int_0^h u_0(z)Z_n(z)\,\dd z .
 \label{eq:ddexp}
\end{align}
Los senos son ortogonales y completos para las condiciones DD; por eso no se
está reteniendo sólo el modo fundamental en~\eqref{eq:ddexp}.
\src{espectro_dd}

\subsection{No deslizante--libre de tensión (DN)}

Con \(Z(0)=0\) y \(Z'(h)=0\), se tiene \(Z=A\sin(qz)\) y
\(\cos(qh)=0\). Por lo tanto,
\begin{align}
 q_m&=\frac{(m+1/2)\pi}{h},&
 Z_m(z)&=\sin\!\left(\frac{(m+1/2)\pi z}{h}\right),\label{eq:dnmodes}\\
 \lambda_m^{DN}&=\nu\left(\frac{(m+1/2)\pi}{h}\right)^2,
 &m&=0,1,\ldots,\label{eq:dnspec}\\
 u(z,t)&=\sum_{m=0}^{\infty}b_m Z_m(z)e^{-\lambda_m^{DN}t},&
 b_m&=\frac2h\int_0^h u_0(z)Z_m(z)\,\dd z .
 \label{eq:dnexp}
\end{align}
También aquí la familia de senos de media onda es ortogonal y completa para DN.
En particular, la condición libre cambia todos los números de onda, no sólo la
tasa más lenta.
\src{espectro_dn}

Como control numérico, para \(\nu=h=1\) las primeras cuatro tasas son
\begin{center}
\begin{tabular}{c@{\quad}rrrr}
\toprule
bordes & \(\lambda_0\) o \(\lambda_1\) & segunda & tercera & cuarta\\
\midrule
DD & 9.86960440 & 39.47841760 & 88.82643961 & 157.91367042\\
DN & 2.46740110 & 22.20660990 & 61.68502751 & 120.90265391\\
\bottomrule
\end{tabular}
\end{center}
Estos decimales se generan con \texttt{python out/checks.py --values} y se
contrastaron contra autovalores de una matriz de diferencias finitas independiente.
\src{valores_espectro}

\section{Reducción bidimensional y fricción de fondo}
\label{sec:reduccion}

Partimos de la ecuación horizontal de Navier--Stokes, con fuerza por unidad de
masa \(\boldsymbol f\),
\begin{equation}
 \partial_t\boldsymbol u_h
 +(\boldsymbol u_h\!\cdot\!\nabla_h)\boldsymbol u_h
 +w\partial_z\boldsymbol u_h
 =-\rho^{-1}\nabla_h p
 +\nu\nabla_h^2\boldsymbol u_h+\nu\partial_{zz}\boldsymbol u_h
 +\boldsymbol f .
 \label{eq:NS3d}
\end{equation}
Al promediar, el término vertical viscoso es exactamente
\begin{equation}
 \frac\nu h\int_0^h\partial_{zz}\boldsymbol u_h\,\dd z
 =\frac\nu h\,[\partial_z\boldsymbol u_h]_{0}^{h}.
 \label{eq:viscprom}
\end{equation}
Hasta este punto no se ha supuesto un perfil vertical.
\src{promedio_ns}

Ahora se hace la clausura explícita de un solo modo,
\begin{equation}
 \boldsymbol u_h(x,y,z,t)=F(z)\,\boldsymbol U(x,y,t),\qquad
 \avg F=1,\qquad \boldsymbol U=\avg{\boldsymbol u_h},\qquad
 \nabla_h\!\cdot\!\boldsymbol U=0.
 \label{eq:ansatz}
\end{equation}
La impermeabilidad y la incomprensibilidad dan entonces \(w=0\). Los perfiles
fundamentales normalizados por su promedio son
\begin{equation}
 F_{DD}(z)=\frac\pi2\sin\!\frac{\pi z}{h},\qquad
 F_{DN}(z)=\frac\pi2\sin\!\frac{\pi z}{2h},\qquad
 \avg{F_{DD}}=\avg{F_{DN}}=1.
 \label{eq:perfiles}
\end{equation}
Para ambos, \(\beta=\avg{F^2}=\pi^2/8=1.233700550136\). La ecuación promediada
queda, bajo esta clausura,
\begin{equation}
 \partial_t\boldsymbol U+\beta(\boldsymbol U\!\cdot\!\nabla_h)\boldsymbol U
 =-\rho^{-1}\nabla_h\avg p+\nu\nabla_h^2\boldsymbol U
 -\alpha\boldsymbol U+\avg{\boldsymbol f}.
 \label{eq:NS2d}
\end{equation}
El factor \(\beta\) aparece porque se promedia \(F^2\), y no se debe ocultar si
la variable 2D elegida es el promedio vertical.
\src{reduccion_modal}

Usando~\eqref{eq:viscprom} y las derivadas de~\eqref{eq:perfiles}, o de manera
equivalente \(F''=-(\lambda_0/\nu)F\), resulta
\begin{align}
 \alpha_{DN}&=\frac{\pi^2\nu}{4h^2}
 =2.467401100272\,\frac\nu{h^2},
 & \alpha_{DD}&=\frac{\pi^2\nu}{h^2}
 =9.869604401089\,\frac\nu{h^2},
 \label{eq:alfas}\\
 \frac{\alpha_{DD}}{\alpha_{DN}}&=4.
 \label{eq:factor4}
\end{align}
Así aparece el término \(-\alpha\boldsymbol U\): es exactamente la acción de la
viscosidad vertical sobre el modo lento supuesto. El valor DD no debe trasladarse
a una capa con tope libre; sobrestima esta predicción por un factor cuatro.
\src{alfas}

Para \(h=6\times10^{-3}\,\mathrm m\), la predicción DN es
\begin{equation}
 \alpha_{DN}=6.8538919452\times10^4\,\nu\ \mathrm{m^{-2}}.
 \label{eq:h6}
\end{equation}
Hace falta insertar la viscosidad cinemática \emph{medida} del electrolito. Sólo
como ejemplo dimensional, si \(\nu=10^{-6}\,\mathrm{m^2\,s^{-1}}\), se obtiene
\(\alpha=0.06853891945\,\mathrm{s^{-1}}\) y \(1/\alpha=14.59025044\,\mathrm s\);
estos dos últimos números no son una medición del experimento.
\src{ejemplo_h6}

\section{Sesgo del PIV de superficie}

En DN, \(Z_0(h)=1\) y
\(\avg{Z_0}=h^{-1}\int_0^h\sin(\pi z/2h)\,\dd z=2/\pi\). Por tanto,
\begin{equation}
 \boxed{\frac{u(h)}{\avg u}=\frac\pi2=1.570796326795},
 \qquad \avg u=\frac2\pi u(h).
 \label{eq:piv}
\end{equation}
Una partícula que flota en el tope mide, en este modelo, una velocidad un
\(57.0796326795\%\) mayor que el promedio vertical. Si se escribe la ecuación en
términos de \(\boldsymbol U_s=\boldsymbol u_h(h)=(\pi/2)\boldsymbol U\), el término
lineal conserva \(\alpha\), mientras el coeficiente no lineal pasa de
\(\pi^2/8\) a \(\beta_s=\pi/4=0.785398163397\).
\src{sesgo_piv}

\section{Cuándo es válida la clausura}
\label{sec:validez}

Sean \(\ell\) la menor escala horizontal que se desea describir, \(U\) una
velocidad típica, \(\tau_H=\ell/U\) el tiempo de vuelco, y
\(\tau_\nu=h^2/\nu\) el tiempo de difusión a través del espesor. Definimos
\begin{equation}
 \epsilon=\frac h\ell,\qquad
 \Rey_\ell=\frac{U\ell}{\nu},\qquad
 \Rey_h=\frac{Uh}{\nu}=\epsilon\Rey_\ell .
 \label{eq:escalas}
\end{equation}
Una condición suficiente para que difusión vertical y geometría delgada dominen
la creación inercial de estructura vertical es el conjunto explícito
\begin{equation}
 \boxed{\epsilon\ll1},\qquad
 \boxed{\frac{\tau_\nu}{\tau_H}
 =\Rey_\ell\epsilon^2=\Rey_h\epsilon\ll1}.
 \label{eq:desigualdades}
\end{equation}
No es necesario imponer \(\Rey_h\ll1\) por separado: el cociente relevante al
balance entre inercia horizontal y viscosidad vertical es
\(\Rey_h\epsilon\). Si además se pretende turbulencia horizontal de alto Reynolds,
la ventana asintótica compatible es \(1\ll\Rey_\ell\ll\epsilon^{-2}\).
Estas desigualdades deben evaluarse con la escala horizontal más pequeña que el
modelo quiera conservar, no sólo con el tamaño de la caja.
\src{escalamiento_validez}

La dominancia del modo lento requiere además que los modos superiores hayan
perdido memoria. Para DN, si \(b_0\ne0\), la condición exacta sobre cada cociente es
\begin{equation}
 \left|\frac{b_m}{b_0}\right|
 e^{-(\lambda_m^{DN}-\lambda_0^{DN})t}\ll1,\qquad m\ge1.
 \label{eq:dominio_modal}
\end{equation}
Para coeficientes iniciales de orden comparable, condiciones suficientes son
\begin{align}
 t&\gg(\lambda_1^{DN}-\lambda_0^{DN})^{-1}
   =\frac{h^2}{2\pi^2\nu} &&\text{(DN)},\label{eq:gapdn}\\
 t&\gg(\lambda_2^{DD}-\lambda_1^{DD})^{-1}
   =\frac{h^2}{3\pi^2\nu} &&\text{(DD)}.\label{eq:gapdd}
\end{align}
Las brechas son, respectivamente,
\(19.73920880218\,\nu/h^2=8\alpha_{DN}\) y
\(29.60881320327\,\nu/h^2=3\alpha_{DD}\).
Para excitación dependiente del tiempo con frecuencia característica \(\omega\),
se requiere también
\begin{equation}
 \frac{\omega}{\lambda_1^{DN}-\lambda_0^{DN}}\ll1
 \quad\text{(o el análogo DD)},
 \label{eq:forcinglento}
\end{equation}
y que la fuerza no inyecte una proyección grande y persistente sobre modos
verticales superiores. Una fuerza estacionaria de perfil arbitrario puede producir
un perfil de Poiseuille o una combinación modal, no necesariamente el modo lento.
\src{validez_modal}

En síntesis: \(\alpha\) es calculable sin ajuste sólo después de fijar las
condiciones de borde, \(h\), \(\nu\) y una clausura vertical válida. La dependencia
\(\alpha\propto\nu h^{-2}\) implica, a primer orden,
\(\delta\alpha/\alpha=\delta\nu/\nu-2\,\delta h/h\); por ello el espesor efectivo
y su incertidumbre importan directamente.
\src{sensibilidad_alpha}

\section{Laplace con Dirichlet abajo y Neumann arriba}
\label{sec:laplace}

Para el modo horizontal
\(e^{\ii(k_xx+k_yy)}\), sea \(k=(k_x^2+k_y^2)^{1/2}\ge0\). La ecuación de
Laplace homogénea se reduce a
\begin{equation}
 \phi''-k^2\phi=0,\qquad \phi(0)=g_0,\qquad \phi'(L_z)=g_1,
 \label{eq:lapode}
\end{equation}
donde \(g_0,g_1\) pueden ser complejos y la prima significa \(+z\).
\src{laplace_reduccion}

Para \(k>0\) se adopta exactamente la base acotada ya usada por
\texttt{laplace\_z}:
\begin{equation}
 \phi(z)=C_1e^{k(z-L_z)}+C_2e^{-kz},\qquad
 \phi'(z)=k\left[C_1e^{k(z-L_z)}-C_2e^{-kz}\right].
 \label{eq:baseacotada}
\end{equation}
Con \(r=e^{-kL_z}\) y \(q=r^2=e^{-2kL_z}\), las dos condiciones forman
\begin{equation}
 rC_1+C_2=g_0,\qquad k(C_1-rC_2)=g_1,
 \label{eq:sistemaDN}
\end{equation}
y su solución es
\begin{equation}
 \boxed{C_1=\frac{g_1/k+r g_0}{1+q}},\qquad
 \boxed{C_2=\frac{g_0-r g_1/k}{1+q}}.
 \label{eq:coefDN}
\end{equation}
Ésta es la normalización solicitada por \(e^{-2kL_z}\): en el intervalo físico,
ambas exponenciales de~\eqref{eq:baseacotada} están entre cero y uno. No se forma
\(e^{+kL_z}\), y el límite \(kL_z\) grande sólo hace \(r,q\to0\), sin desborde.
La suma del denominador, en vez de la resta de la rama DD, es consecuencia de la
condición Neumann superior.
\src{coeficientes_dn}

El modo horizontal constante se debe separar antes de dividir por \(k\):
\begin{equation}
 k=0:\qquad \phi(z)=g_0+g_1z,\qquad \phi'(z)=g_1.
 \label{eq:k0}
\end{equation}
En la rama lineal de SPECTER esto corresponde a
\(\texttt{coef(1,j,i)}=g_1\) (pendiente) y
\(\texttt{coef(2,j,i)}=g_0\) (intercepto).
Además, aunque~\eqref{eq:coefDN} es estable para \(kL_z\) grande, \(C_1\) y
\(C_2\) se cancelan cuando \(k\to0\). La implementación evalúa los modos pequeños
mediante la forma regular equivalente
\begin{equation}
 \phi(z)=g_0\{\cosh(kz)-\tanh(kL_z)\sinh(kz)\}
 +g_1\frac{\sinh(kz)}{k\cosh(kL_z)},
 \label{eq:formaregular}
\end{equation}
calculando \(\sinh(kz)/k=z\,\mathrm{sinhc}(kz)\). Así se cubren sin cancelación
tanto \(k=0\) como valores positivos muy pequeños, mientras la base exponencial
acotada cubre los valores grandes.
\src{limite_k0}

La correspondencia para una futura rama Fortran es
\(\texttt{coef(1,j,i)}=C_1\),
\(\texttt{coef(2,j,i)}=C_2\),
\(\texttt{bc(1,j,i)}=g_0\) y
\(\texttt{bc(2,j,i)}=g_1\). En esta fase no se modificó el archivo Fortran.
\src{convencion_specter}

\section{Implementación y verificaciones}

\texttt{out/capa.py} expone exactamente
\texttt{tasas\_decaimiento},
\texttt{cociente\_superficie\_promedio}, \texttt{alfa},
\texttt{solver\_1d} y \texttt{coef\_laplace\_dn}. Sólo acepta las cadenas de
borde congeladas \texttt{noslip-noslip} y \texttt{noslip-libre}.
\src{api_capa}

El solver usa \(n_z\) intervalos uniformes, diferencias centradas de orden dos,
el nodo fantasma \(u_{n_z+1}=u_{n_z-1}\) para DN, y aplica
\(u(t)=\exp(-At)u(0)\) al sistema semidiscreto. Por tanto sí integra un operador
FD; no devuelve la solución continua insertada a mano. Al duplicar resolución de
48 a 96 intervalos, los errores de tasa fueron
\(3.523\times10^{-3}\to8.808\times10^{-4}\) (DD) y
\(2.202\times10^{-4}\to5.505\times10^{-5}\) (DN), con orden medido 2.000.
\src{solver_fd}

Los chequeos propios pasaron 5/5. En particular, siete casos Laplace DN ---incluidos
\(kL_z=700\), \(k=10^{-10}\) y datos complejos--- dieron residuo de borde máximo
\(4.441\times10^{-16}\). La puerta de aceptación fijada por hash pasó 6/6; sus
comparaciones de espectro tuvieron errores relativos máximos
\(8.23\times10^{-5}\) (DD) y \(6.30\times10^{-5}\) (DN), y su solver midió orden
2.00. Todas las ejecuciones fueron secuenciales.
\src{checks_propios,puerta_aceptacion}

\section{Procedencia de los pasos}

Se \textbf{derivaron en este trabajo}: las dos condiciones seculares, ambos
espectros y coeficientes de expansión; el promedio viscoso y \(\alpha\); el factor
PIV y \(\beta\); las desigualdades por escalamiento; y el sistema algebraico DN,
sus coeficientes acotados y su límite \(k=0\). Los números impresos provienen de
\texttt{out/checks.py::derived\_values} o de ejecuciones de checks registradas.
\src{procedencia_derivada}

Se \textbf{recordaron de literatura estándar}, y luego se rederivaron aquí para no
importar un coeficiente sin prueba, el método de separación de variables y la
completitud Sturm--Liouville para la ecuación de difusión, así como las ecuaciones
de Navier--Stokes, el promedio vertical y la interpretación de una superficie sin
tensión tangencial. De SPECTER se tomó sólo la convención computacional de la base
exponencial acotada y el orden de \texttt{coef}; no se tomó de allí la rama DN,
porque no existe en el código inspeccionado.
\src{literatura_estandar,convencion_specter}

\section{Salvedades y cosas no verificadas}
\label{sec:salvedades}

\begin{itemize}
\item No se accedió a datos experimentales ni se midieron \(\nu\), el espesor
  efectivo, la rugosidad o el posible deslizamiento. Por eso~\eqref{eq:h6} es una
  predicción paramétrica, no una calibración de la celda.
\item No se verificaron deformación superficial, ondas, evaporación, surfactantes,
  gradientes de tensión superficial (Marangoni), viento, ni una tapa parcialmente
  contaminada. Cualquiera de ellos cambia la condición \(\partial_z\boldsymbol u_h=0\).
\item El ansatz de un modo no es una identidad para Navier--Stokes no lineal ni para
  forzamiento vertical arbitrario. No se hicieron DNS 3D; las condiciones de la
  sección~\ref{sec:validez} son condiciones asintóticas suficientes, no una frontera
  universal medida para la transición a flujo cuasi--2D.
\item Se usó viscosidad molecular constante y una capa de profundidad uniforme. No
  se modelaron una capa límite turbulenta, viscosidad de remolino, estratificación,
  dos líquidos superpuestos ni paredes laterales.
\item La rama Laplace DN se verificó como módulo Python por residuos de borde,
  finitud y límite pequeño. No se implementó ni compiló en Fortran, no se ejecutó
  MPI y no se verificaron todavía distribución espectral, índices ni precisión
  \texttt{GP} dentro de SPECTER.
\item La convergencia FD y la puerta numérica no constituyen una prueba formal para
  todo dato o parámetro. Sí prueban los casos enumerados y complementan las
  sustituciones algebraicas de esta nota.
\end{itemize}
\src{salvedades}

\begin{thebibliography}{9}
\bibitem{carslaw}
H. S. Carslaw y J. C. Jaeger,
\emph{Conduction of Heat in Solids}, 2.a ed., Oxford University Press, 1959.

\bibitem{kundu}
P. K. Kundu, I. M. Cohen y D. R. Dowling,
\emph{Fluid Mechanics}, 6.a ed., Academic Press, 2016.

\bibitem{specter}
SPECTER upstream, \texttt{src/boundary/boundary\_mod.fpp}, subrutina
\texttt{laplace\_z}, inspeccionada en modo sólo lectura.
\end{thebibliography}

\end{document}
